\documentclass[12pt]{article}

\usepackage[a4paper, total={6in, 8in}]{geometry}
\usepackage{textcomp}

\begin{document}
\setlength{\parindent}{0pt}

\def \t {\textrightarrow}
\def \v {\vspace{1em}}
\def \bi {\begin{itemize}}
\def \ei {\end{itemize}}
\def \s[#1] {\section*{#1}}
\def \ss[#1] {\subsection*{#1}}
\def \sss[#1] {\subsubsection*{#1}}

\s[Ripasso]
\ss[Eta giolittiana]
Siamo all'inizio del 900 (ma primo governo è tra i due governi Cripsi) \t sinistra storica ma ha tratti peculiari \\
Riprendere fasci siciliani, Crispi era molto autoritario anche nei confronti dei fasci \\
Anni tra il 1903 e 1914, ma era gia ministro nel governo Zanardelli nel 1901 \\
Non sono governi continuativi \t uno dei tratti distintivi è il dimettersi: quanco capisce di non essere gradito dall'opposizione e sue proposte vengono cassate perchè legate al suo nome, lui mette al suo posto uno manovrabile \t e poi torna \\
Altri caratteri:
\bi
  \item dimissioni
  \item doppio volto: nord riformista, progressista al sud prefetti che hanno controllo e disinteresse (Salvemini dice ministro della malavita)
  \item politica del non internvento \t non seda situazioni difficili sociali, perchè la cosa si risolve da sola \t con i fasci siciliani non sara cosi, 1904 primo grande sciopero italiano funzione
\ei
Sbaglia quando nel 1920, richiamato al governo, sottovaluta pericolo del fascismo e in un confronto tra stato e operai favorisce la classe operaia \t questo è errore: imprenditoria si sposta verso le destre, e Mussolini quini \\
Infine trasformismo: trovare alleanze parlamentari per continuare ad avere la maggioranza \t trasformismo viene inaguruato da De Pretis, ed è prassi della sinistra storica \t ma Giolitti opera così \\
Nel 1903 chiede a Turati (socialista leader) di entrare nel suo governo \t mossa politica, Turati non poteva accettare o si rompeva il partito \\
Anche patto Gentiloni, ovvero accordo del 1913 con unione cattolica ?? che permette ai cattolici di votare liste liberali in funzione anti comunista \\
Cosi Giolitti si rafforza \t pero poi si dimette, mette Salandra, ma poi non riesce a tornare \\
Trasformismo in politica estera invece rimane in triplice alleanza, pero mantiene contatti con potenze dell'intesa \t questo permettera il patto di Londra \\
Giolitti cerchera riscatto con la guerra di Libia, che si inserisce in un contesto imperialista (col disfacimento dell impero ottomano) \t guerra di Libia nel 1912 \\
In questi anni i balcani si sta ribellando, quindi impero ottomano ha due fronti \\
Ultimo govenro giolitti è già dentro questa eta, che è quella dell'imperialismo \t momento che genera le cause remote della WWI (omicidio di Sarajevo, casus belli ma non la causa) \\
Giolitti concede anche il suffragio universale maschile, in funzione antisocialista \t aveva bisogno di togliere al socialismo il punto di forza del partito, che puntava sul suffragio per avere consenso \\
Il partito socialista è diviso in massimalisti e riformisti (nel 75 partito dei lavoratori, il partito socialista nasce nel 95): massimalisti vogliono rivoluzione \\
Si alternano alla guida del partito \t quando Turati guida, GIolitti ha interlocutore \t se no no \t nel 1912 partito massimalista ha controllo, quindi suffragio per colpirlo \\
Patto Gentiloni vince, ma poi si dimette \t va al suo posto Salandra che non è il burattino di Giolitti \t affronta la settimana rossa, e poi WWI, che si ineserisce nel quadro dell'imperialismo

\ss[Imperialismo]
Congresso di Berlino è del 78 è importante per i Balacani \t voluti da Bismark, voleva risolvere questioni in Africa e Balcani \t al congresso di Berlino Austria ha protettorato sulla Bosnia (che poi annette) \\
Nazionalismo, Germania di Guglielmo II, delinea un quadro che culmina con l'omicidio si Sarajevo (casus belli, non causa) \\
Cause sono in sitauzione internazionale esplosiva che trova radici nell'imperialismo

\ss[Prima Guerra Mondiale]
Omicidio di Sarajevo quando? \t ultimatum inaccettabile dell austria alla Serbia \t vuole commissione mista di indagine sull omicidio di Sarajevo \t era fatto per non essere accettato (48 ore) \t serbia rifiuta, perche è appoggiata dalla russia \\
La russia aveva gia appoggiato la lega balcanica \t sosteneva indipendenza di alcuni stati, tra cui la serbia \\
Austri dichiara la guerra alla serbia, gioco alleanze triplice alleanza (dell 1882) e triplice intesa (del 1907) \t strano che russia è in intesa, perchè è impero zarista tra paesi democratici \t ma in realtà russia è grande nemica dell austria \\
Quindi scatta con la Germania, che si schiera (italia no perchè l'alleanza è difensiva) \t mentre poi c'è Russia con Francia etc. \\
Italia sta fuori, dibattito interventismo vs. neutralismo \t alcuni sono a favore con la triplice alleanza altri anche con intesa \\
Triplice intesa promette territori a Italia, mentre proposte della alleanza non sono forti \t col patto di Londra nel 1915 Italia entra in guerra con intesa \\
Primi 6 mesi guerra di movimento \t poi guerra di trincea \t guerra che inizialmente a favore degli imperi centrali, ma poi nel 1917:
\bi
  \item entra in guerra USA \t guerra marina tra germ e uk, che fa esplodere Lusitania
  \item rivoluzione russa \t Brest-Litovsk firmato armistizio, ma gia quando Lenin va al potere ritira
  \item caporetto \t truppe austriache vanno quindi in italia, non piu problema russo (in realta prima 1916 spedizione punitiva)
\ei
Disfatta di caporetto penetra per 150 km in territorio \\
DOpo Lusitania Wilson decide di entrare in guerra (anche se prima forte isolazionismo) \\
1917 cambia sorti della guerra \t guerra si conclude con battaglia nel Vittorio Veneto e battaglia di Amien \\
Russia sconfitta ai Laghi masuri nel 15, quindi crolla serbia perche entra anche bulgaria con austria \t battaglia di Verdan e Some del 16 \\
I tedeschi lanciano offensiva a Verdan, mentre francia risponde con somme \t entrambe battaglie fanno strage \\
Poi battaglia di caporetto, poi vittorio veneto e Amien \t Retonde e Villa giusti

\ss[Trattati di pace]
14 punti di Wilson, che includono 2 punti: 
\bi
  \item autodeterminazione dei popoli, smetito dai trattati di pace
  \item societa delle nazioni, deve risolvere diplomaticamente (nasce nel 1919, ma USA rimane fuori) \t non funziona perche la nazione che la propone, non partecipa
\ei
Società delle nazioni non funziona anche perche non ha apparato militare \t solo sanzioni come in guerra etiopia \\
Trattato conclusivo è Versaille \t punisce la germania e considerata unica responsabile \t ricordare maggioranza delle clausole, soprattuto privazioni alla Germania:
\bi
  \item 130 miliardi di franchi d oro come riparazione
  \item smilitarizzazione
  \item occupazione della Saar e delle miniere della francia
  \item colonie perse
  \item perde Alsazia e Lorena a Francia
  \item Danzica alla polonia
  \item vietata unificazione con l'Austria (paura di possibile unificazione e creazione di un centro forte in europa)
\ei
Trattato di Versaille non riguarda l'italia \t poi ci sono altri trattati, noi firmiamo quello di Saint Germain \t permette di annettere territori irredenti previsti dal patto di londra, ma non tutti gli altri territori \t se no si andava contro autodeterminazione popoli \\
Cosi Vittorio Emanuele Orlando forza prima la situazione chiedendo Fiume, quindi abbandona trattative \t mentre italia è assente c'è sparttizione colonie tedesche \t poi torna e firma trattato \\
Tutto questo in esaltazione nazionalistica e vittoria mutilata \t quindi occupata città di Fiume, e D'Annunzio guida nazionalisti a occupare settembre del 19 a occupare Fiume (come atto di protesta) \\
Città di Fiume presa, istituito consiglio nazionale \t si chiede annessione all'italia di Fiume, ma viene respinta \t viene anzi imposta la liberazione di Fiume \\
Inoltre grande tensione tra Croati e Italiani, perchè era nato il paese della Croazia \\
Nessuno oppone resistenza, e questo genera situazione difficile per l'italia \t a livello internazionale appariva che l'italia non aveva accettato condizioni, e si era spinta troppo oltre \\
Al governo c'è Nitti (infatti Orlando si era dimesso) \t economista, non riesce a gestire tensione internazionale \t Mussolini critica Nitti, che decide di intervenire inviando truppe alla frontiera guidate da Badoglio (che non è d'accordo) \\
Infatti c'erano molti volontari (e soldati) che andavano a fiume, e questo poteva generare ribellione nell'esercito \\
Nitti arriva alla frontiera, mentre D'Annunzio non riesce a fare azione politica concreta \t arringa le folle, ma non riesce a fare molto \t stesa la costituzione che è Carta del Carnaro (molti visionaria, e anche moderna: parla dell'uguaglianza tra uomo e donna per esempio) \\
Cade governo Nitti, Giolitti chiamato e nel 20 viene firmato trattato di Rapallo
\bi
  \item fiume città libera
  \item Jugoslavia ottiene dalmazia 
  \item italia ottiene istria
\ei
Dopo alcuni scontri, Fiume liberata \t Mussolini dice che è azione fraticidia questa, pero dice anche che Rapallo era unica soluzione \\
D'annunzio nel 20 scrive un documento "Programma d'annunziano per una insurrezzoine e marcia su roma" \t descrizione del suo progetto, Fiume era un primo step per poi attuare colpo di stato sotto la sua guida \t a Mussolini non piace, lo vuole fare lui \\
C'è simpatia, ma anche distacco di Mussolini \t lo usa per propaganda, ma poteva diventare un concorrente

\ss[Rivoluzione russa]
Motlo arretrata, fortemente agricolo \t si tentano riforme a partire dal 1861 con Alessandro II \t abolisce servitu della gleba senza passaggi intermedi, riforme falliscono \\
Inizia sviluppo industriale solo in alcune zone (le grandi città, Mosca San pietro burgo e Baku) \t stato fortemente accentrato, si crea opposizione tra occidentalisti e slavofili \\
\bi
  \item occi. vogliono capitalismo, democrazia etc. 
  \item slavofili \t vogliono segure risorse interne, quindi agricoltura \t movimento populista che mette al centro i contadini ma fallisce
\ei
Movimento sara fallimentare e ci saranno pero anche frange terroristiche (populista uccide lo zar) \t sono i social rivoluzionari \\
Opposizione marxista esiliata, che punta alla classe operaia (non ai contadini) \t nel 1898 nasce il partito operaio social democratico che dal 1903 si divide in due correnti contrapposte:
\bi
  \item bolscevichi \t lenin
  \item menscevichi \t martov
\ei
Nel "Che fare" del 1903 \t Lenin dice che:
\bi
  \item rivoluzione ha priorita su rivendicazioni sindacali
  \item rivoluzione non è spontanea, servono conduttori
  \item partito deve essere monolitico
\ei
Menscevichi vogliono invece la democrazia tedesca, riforme etc. \\
A partire dal 1903 diventeranno due partiti distinti \\
Il fondatore del comunismo russo è pero Plekanov (in esilio si fa portavoce del discorso marxista, e dice che è necessaria rivoluzione capitalista prima di quella comunista \t in russia non c'è propletariato perche è ancora indietro) \\
La russia si presenta composita dal punto di vista delle opposizioni \t e ci saranno 3 riv.:

\sss[1905]
Rivoluzione spontanea che fallisce \t scoppia per mancanza di cibo a San P. \t processione che arriva fino al palazzo d'inverno, aperto fuoco sulla folla: "Domenica di sangue" \t questo crea rivolte e scioperi \t opposizione sale anche nella borghesia, nascono i cadetti che vogliono sistema costituzionale moderato \t lo zar si spaventa \\
Permette la Duma \t assemblea consultiva che pero non avra alcun valore, ma accontenta i borghesi (che comunque non vogliono riv. socialista) \\
Protesta si allarga, corazzata Potionki si ammutina \t a San P. si crea il primo soviet = consiglio dei lavoratori, con capo Trozky \\
Soviet è per rivendicazioni sindacali, ma anche organo politico \t col soviet la rivoluzione ha cambiato aspetto: prima era di rivendicazione, adesso è anche politica \\
Tra 06 e 17 elette diverse Dume, tutte sciolte dallo Zar quando si contrappongono \t poi russia va in guerra inpreparata militarmente ed economicamente \t diminuisce produzione del grado, etc. \\
Nel 1915 la russia perde il controllo dei territori polacchi, nuovi scioperi \t la riv. del 05 nasce anche dalla guerra russo-giapponese che genera crisi

\sss[1907]
Nel febbraio c'è prima rivoluzione \t sempre spontanea, esplode perche insorgono gli operai \t si chiede la distribuzione della guerra e democrazia \\
Lo zar non riesce a controllare, abdica nel marzo e quindo governo provvisorio democratico \t nasce repubblica con principe Lvov, ma a fianco del governo c'è il soviet \\
Tra febbraio e ottobre c'è dualismo di potere, che va verso il crollo della duma e una crescita del soviet \t l'istituzione credibile per il popolo è il soviet
\bi
  \item governo vuole continuare la guerra perche vittoria rafforza stato \t dopo guerra voleva vaga politica di riforme
  \item soviet vuole sconfiggere potenze imperialiste come austria e germ. \t per difendere la riv. \t soviet punta a riforma agraria dopo guerra
\ei
Aprile del 17 Lenin torna dalla svizzera in russia \t scrive le "Tesi di aprile" = documento operativo, spiega cosa bisogna fare:
\bi
  \item dare potere ai soviet = distruggere governo provvisorio
  \item uscire dalla guerra 
  \item riforma agraria = terre ai contadini
  \item riv. non è ancora pronta, perche senza coinvolgimento del mondo contadino riv. fallisce \t riv. deve essere proletaria, ma maggioranza contadini
\ei
Riv. viene preparata \t anche grazie all andamento della guerra: gov. provvisorio vuole fare offensiva contro austro-tedeschi che fallisce, truppe russe vanno all assalto ma macello \\
Popolazione si oppone a far partire soldati dopo questo evento \t poi Cornilov tenta colpo di stato, ma il governo presieduto da Kierensky previene colpo di stato con operai, contadini e bolscevichi (che cosi conquistano maggioranza nei soviet di mosca e pietro grado) \\
Si crea la Guardia Rossa da Trozky \t si preparano piani per la rivoluzione, e il 24 ottobre senza spargimento di sangue bolscevichi ottengono controllo dei palazzi di governo \\
Al congresso panrusso viene procalamato il potere sovietico, il potere dei soviet \\
Lenin approva:
\bi
  \item decreto sulla pace: immediato armistizio, si rinuncia a tutto
  \item decreto sulla terra: abolizione proprietà privata
\ei
Consiglio dei commissari del popolo con Lenin, nazionalizzate la banche e i mezzi di produzione agli operai per poi nazionalizzarli \\
Il 17 si tengono elezioni per assemblea costituente \t campagne si schierano con i social rivoluzionari, sfavorevole per bolscevichi \t quindi bolscevichi la sciolgono (giustificata da Lenin con dittatura del proletariato) \\
Si legge tentativo autoritario \t molti scappano dalla russia \\
Nel 18 pace di Brest-Litovsk \t condizioni durissime:
\bi
\item alla germania regioni tra bielo e caucaso
\item ucraina e finlandia indipendenti
\item russia rinucnia a pretese terr. sui baltici e polonia
\ei
Russia rinuncia quindi a territori fertili e ricchi \\
Ai social rivoluzionari la piace non piace, escono dal governo \t dopo pace i bolscevichi sono gli unici a governare \\
Dal 18:
\bi
  \item guerra civile: zar con armata bianca e armata rossa \t 3 milioni di morti, tra cui Zar e tutta sua famiglia (Nicola II)
  \item potenze occ. minacciano russia, vogliono far cadere Lenin \t potere socialista qua è pericoloso \t alti e bassi, ma Lenin vince grazie all'appoggio dei contadini (Lenin è il male minore piuttosto che regime zarista)
\ei
Nel 20 Polonia cerca di riappropiarsi dei territori che aveva perso con Versaille \t armata rossa respinge e arriva fino a Varsavia, guerra si conclude nel 21 e russia riacquisice alcuni territori \\
Nel frattempo prima cost. sovietica:
\bi
  \item russia è repubblica federale
  \item si possono aggregare altre repubblica negli altri territori
  \item 1922 nasce URSS
\ei
Potere bolscevico diventa sempre + autoritario \t creata la CECA = poliza politica che diventa famosa per la sua violenza \\
Fin'ora applicato il comunismo di guerra: 
\bi
  \item industrie nel caos, operai hanno pieno controllo 
  \item governo non riesce a riscuotere tasse, quindi stampa = inflazione
  \item contadini producevano per la loro soppravvivenza, non davano alle citta
\ei
Nel 21 inaugura la NEP:
\bi
  \item contadini possono coltivare terra per loro necessita, dare allo stato e vendere eccedenza
  \item legalizzato mercato nero
  \item stato mantiene il controllo delle fabbriche grandi, per quelle piccole rientrano come soci i vecchi propretari (per questo verra attaccata, nep reintroduce forma di proprieta privata)
\ei
Nep inoltre è a vantaggio solo dei contadini quindi criticata \\
Nep funziona pero, e nel 1926 la produzione agricola e industriale tornano ai livelli del 1914 \\
Viene proibito nel frattempo frazionismo = no corrente stabili nel partito \t la posizione del centralismo democratico (discussione avviene nel partito, ma quando partito assume posizione diventa vincolante per tutti, anche per chi non la condivideva) \t condivisa anche da tutti i partiti comunisti europei \\
Questo vuol dire fine democrazia interna, e accentuare carattere autoritario \t urss va verso stato autoritario a partito unico, il CUS \\
Nel 22 Lenin viene colpito da un ictus e muore nel 24 \t si apre conflitto tra Trozky (ideologo) e Stalin (segretario) \t si scontrano su
\bi
  \item conduzione partito (democratico vs autoritario)
  \item nep (non sostiene, altro si)
  \item esportazione della rivoluzione (stalin no, vuole socialismo in un solo paese)
\ei
Linea di stalin prevale \t lui aveva dialogato con i paesi \\
Mantenuta la nep fino al 27 \t stalin fa:
\bi
  \item industrialiazzazione collettivizzazione forzata \t mette in atto dopo le nep, prende le idee di trozky
  \item culto del capo \t forte propaganda su di lui, fa riscivere i libri di storia
  \item eliminazione dell'opposzione \t che pero danneggia molto urss, tra il 25 e il 38 da vita alle purghe che colpiscono tutti oppositori \t sara motivo di debolezza alle soglie della WWII, perche ufficiali sono stati dimezzati
\ei
Stalin vuole rendere urss la prima nazione al mondo \t tutti i settori devono essere sacrificati all'industrialiazzazione, che procede per piani quinquennali \\
Industrialiazzazione significa grandi investimenti, risorse alimentari e agricole subordinate (o comunque indirizzati al decollo industriale) \t comunque il primo (si rivolge all industria pesante) e il secondo hanno grande successo \\
Terra divisa in sovkozi e kozi ??? \t nei primi i contadini usano la terra in modo colletivo ma terra è dello stato, mentre sovkozi erano interamente statali e contadini sono dipendenti dello stato e maggior parte contributi va a stato \\
Stato a fine anni 30 controlla interamente campagne, e non esiste attivita economica indipendente dallo stato \\
Urss diventa una grande potenza industriale \\
1919 c'è la terza internazionale (lenin fonda per ??? vedere, anche internazionale fascista con guerra spagna)

\ss[Italia e Germania dopo WWI]
Si scatena confronto acceso tra destra e sinistra in entrambi \t Biennio Rosso tra 20 e 22 anni che si pensa che si possa creare una rivoluzione in italia \\
Le forze socialiste prendono piede, e c'è clima da guerra civile \t lo stato liberale non riesce a gestire \t nascono consigli operai, crisi sociale \t orizzonte in cui c'è sfiducia nella democrazia liberale, che si ritiene un governo che non funziona + \\
Situazione analoga, lega di Spartaco in Germania che org. meglio \t esiti diversi: in italia l'esito del biennio rosso è il fascismo \\
Infatti giolitti fa errore, dopo Nitti si trova a gestire situazione economica disastrosa, debito pubblico altissimo \t e giolitti media tra tutte queste forze andando ad agevolare richieste degli operai, avvantaggiandoli dal punto di vista legislativo \\
Classe imprenditoriale si sposta su posizioni + conservatrici \\
In germania non succede questo \t il tentativo di colpo di stato della lega di Spartaco viene fermato dalla repubblica di Weimar, che con Stresemann ha stabilita (finche lui non muore) \\
Evoluzione tra i paesi è diversa

\ss[Italia dopo WWI]
1920 occupazione della terra e delle fabbriche \t FIOM chiede rinnovo contratto con operai metalmecc \t richieste rifutate, quindi sciopero bianco: entrnao in fabbrica ma non lavorano \\
Imprenditori chiudono le fabbriche, operai le occupano \t nel nord vengono occupate tantissime, e in alcuni casi tentano di continuare a produrre \t l'idea era che doveva essere l'inizio di una rivoluzione \\
Cosa difficile in italia pero, movimento non è chiaro ideologicamente e strategicamente \t quindi non si riesce a realizzare \\
Ordine nuovo, rivista di Gramsci e Togliatti \t rivista di matrice comunista che indica il percorso \\
A livorno nel 21 massimalisti si staccano e si attaccano a ordinovisti \t mentre abrogato non expedit, Don Luigi Sturzo, chiesa ha paura quindi permette voto \\

\ss[Presa al potere di Mussolini]
Si crea movimento dei fasci italiani di comabattimento \\
Programma di san sepolcro dioc ane fonda movimenot confuso ?? \\
Mussolini va al potere con marcia su roma, che pero è il punto di arrivo di una affermazione territoriale del fascismo in atto da anni \\
Il fascismo aveva preso possesso di molte zone di campagna, ma anche di città con le amministrazioni comunali \t la sua presenza sul territorio è consolidata \\
Nascono subito le squadre fasciste, che riportavano l'ordine nelle campagne, facevano paura e seminavano il terrore \\
Eccidio di Bologna, dopo questo nasce fascismo agrario = le squadre dilagano \t a Bologna vincono socialiti ad amministrative, sindaco saluta sulla piazza, partono colpi di pistola, scappano tutti, servizio d'ordine socialista spara sulla folla \\
Alcuni dicono che primi spari erano fascisti \t da qui azione squadre dilaga, vuole riportare l'ordine \\
Il governo liberale inoltre sottovaluta, pensa di poter usare Mussolini per consolidarsi \t si vuol fare risolvere la situazione a mussolini \\
Giolitti quindi indice nuove elezioni \t nel 21 accetta che vengano cretae delle liste comuni (formate da liberali, centro e fascisti mentre continuano ad usare violenza nelle campagne) \\
Giolitti fa questo perche vuole ridimensionare socialisti e popolari, ostili al liberalismo \t popolari aumentano consensi, mentre fascisti entrano in parlamento \\
Giolitti sperava di acquisire controllo ma non riesce \\
\bi
\item il movimento diventa partiton nazionale fascista
\item bisogna controllare Balbo, Farinacci \t azione + aggressiva delle squadre
\item crea milizia paramilare salute nazionale
\ei
Governo liberale muore:
\bi
\item bonomi governa per 6 mesi 
\item poi Luigi Facta \t appoggiato da liberali, ma governo debole \t le forze di maggioranza non coese e lui non figura decisa
\ei
Facta si trova la marcia su roma \t Mussolini riunisce squadre a napoli, Facta chiede statod i assedio al re (ma rifiuta) \t quindi le squadre con triumvirato entrano a roma nel ??? \\
Mussolini che era pronto a scappare torna \\
Mussolini nel frattempo da repubblicano diventa monarchico, da socialista filocapitalista, abbandona anticlericalismo \t cosi partito fascita è + presentabile politicamente, quindi fa marcia su roma

\ss[Fascismo]
In fase legalitaria si muove nelle strutture dello stato liberale (non fa cose legali pero) \\
Nella fase legalitaria fa riforma elettorale, ovvero legge Acerbo (novembre del 23) \t la lista che ha maggioranza relativa alle elezioni ha 2 terzi seggi in parlamento \\
In seconda battuta la maggioranza relativa verra definita al 25 per cento \\
Alle elezioni pero stra vince quindi Acerbo non serve, anche se è legge iniqua \\
Metteotti, del partito socialista, critica Acerbo, elezioni, accusa mussolini (nel maggio 24) \t viene ucciso \\
Questo è il momento di risveglio del parlamento e della popolazione \t mandante non scoperto, mentre esecutori arrestati \t questo è il momento di crollo del consenso del partito \\
Partito com. vuole proclamare sciopero ma rifiutato \t quindi partiti non vogliono + entrare in parlamento (secessione dell'Amentino che riprende quella dei romani) \\
I seccessionisti sperano che il re intervenga, ma non lo fa \t mussolini chiude parlamento, lo riapre, loro non tornano (lui ha dato possibilita quindi da colpa a loro) \\
Nel 25 fa famoso discorso \t se fascismo è associazione a delinquere etc. \t inizio fase totalitaria \\
Governo totalitario = organi di partito si sovrappongono agli organi dello stato \t partito unico, e diffusa una sola idea a tutti i livelli \\
Dal 25 leggi fascistissime, gran consiglio del fascismo \t cambiano istituzioni stato \t con propaganda, istruzione etc. \\
Economicamente: fino al 25 politica liberista, mentre dal 25 politica protezionista (vuole autarchia, alimentare ma non in funzione della guerra) \t per autarchia fa bonifica delle regioni, battaglia del grano con francia, costruisce citta etc. \\
Prevede anche sistema corporativo \t deve stabilire dialogo tra operai e imprenditori \t ma fallisce \\
Nascono apparati statali che devono aiutare industria:
\bi
  \item Ini \t istituto di credito che sostiene aziende 
  \item Iri stato diventa azionista di maggioranza di molte grandi aziende e banche
\ei
Fascismo è un totalitarismo imperfetto perchè c'è il re \t istituzione monarchica è presente \\
Fascismo non ha vera presa sul territorio nazionale perche alcuni valori cozzano con quelli diffusissimi della chiesa \\
Nel 29 patti laternansi
\bi
\item concordato: stato confessionale
\item trattato: chiude questione romana, iniziata nel 1861 come rapporti mentre nel 70 breccia di porta pia escalation \t tattato tra stati
\item indennizzazione finanziaria: chiesa indennizzazione per territori che perde
\ei

\ss[Politica estera fascista]
Relazioni cordiali, amico di tutti \t fa da mediatore \t nel 25 conferenza di Locarno con germania, quando Hitler fa cose chiamato mussolini (come prima tentata annessione dell'austria) \t relazioni diplomatiche positive \\
Poi guerra d'etiopia \t si incrinano rapporto con francia e ing. \t e avvicinamento con la germania di hitler, unico a sostenerlo quando societa nazioni fa sanzioni (inutili) \\
Guerra d'etiopia, decisa nel 34 ma truppe partono nel 35 \t no dichiarazione di guerra, nel 36 conquestata Adis Abeba \\
Nel 36 annuciata la nascita dell'impero dell africa orientale italiana \t è annessione di successo? no: etiopia non ha risorse utili al momento \\
Ritirate le sanzioni, perche ormai era stata annessa \t francia e ing. riconoscono impero etiopia = mussolini ha vinto \\
Nel 36 nasce anche asse roma-berlino = consolidamento e guerra di spagna \\
Avvicinamento sancito nel 38 con le guerre razziali, che riprendono quelle di Norimberga \t questo genera opposizione in italia (opposizione è pero reato ora) \\
Benedetto Croce in tutto questo invece si scontra con il fascimo \t era intellettuale molto stimato e famoso in europa \t quindi era tollerato dal regime, che non vuole danneggiarsi \\
Croce stampa la rivista "La critica" \t per tutto il ventennio fascista \t fa opposizione teorica, non politica in realta \t bobbio dice che croce è stata la coscienza antifascista in italia \\
Nel 37 asse roma berlino tokio, patto d'acciaio nel 39 (patto scellerato) \t guerre difensiva e offensiva, italia si impegna \\
Con patto d'acciaio si compie la sottomissione dell italia fascista a quella tedesca

\ss[Germania dopo WWI]
Nelle sinistre c'è spaccatura tra componente maggioritaria della social democrazia (democratica), la lega di Spartaco con Rosa Louxemburg (socialismo + intransigente) \\
La lega da vita nel 1918 al partito comunista tedesco \t gli spartachisti vogliono generare rivoluzione (nel gennaio 1919) \t represso aspramente dalle forze dell'ordine, molti spartachisti e Rosa giustiziati \\
Non è come il biennio rosso \t la crisi della germania viene infatti risolta dalla componente moderata \\
Assemblea costituente eletta, si riunisce a Weimar e nasce la repubblica in 19:
\bi
\item 17 regioni
\item parlamento eletto a suffragio univerale proporzionale
\item potere esecutivo al cancelliere, responsabile davanti al parlamento (parlamentarismo)
\ei
Aggiungere punti \\
Sistema è formalmente parlamentare, ma ha caratteri presidenziali molto forti \t che con i suoi poteri puo snaturare l impianto parlamentare \\
Il presidente puo 
\bi
\item nominare il cancelliere
\item comandare esercito
\item indirre referendum
\item aritcolo 48: da poteri straordinari in situazioni emergenze (ma generiche)
\item inoltre 48 dice che presidente in emergenza puo sospendere i diritti fondamentali (Hitler)
\ei
Trattato di Versaille genera motivi di forte mal contento \t 
\bi
\item profonda crisi economica
\item passa da modello monarchico a repubblicano che schricchiola
\item germania vive tensioni rivoluzionarie, che minano stabilita
\ei
Tra il 19 e il 23 ci sono gravi tensioni, date da:
\bi
\item movimento spartachista
\item soprattutto da non accettazione del trattato di Versaille \t governo accusato
\ei
Si tenta un colpo di stato, nel 20 provato da Kapp (estrema destra) \t fermato da forze democratiche \t poi putch di Monaco di Hitler \\
Hitler è del partito ? \t tenta colpo di stato, bloccato sul nascere da Stresemann (grande fautore della stabilita della rep. di Weimar) \\

\ss[Stresemann]
Stresemann rende la germania uno stato anche diplomaticamente accettato \t lui è il leader del partito popolare (democratico-liberale) \t nel 23 fa governo di coalizione con partito centro e socialisti \\
Vuole costruzione finanziaria, sanare situazione sociale = nel complesso dare stabilita
\bi
  \item fa riforma monetaria e sostituisce il franco
  \item francia
  \item reprime con la forze estremismo
  \item viene aiutato da USA, con triangolazione finanziaria (USA, USA Germania, USA UK FR)
\ei
USA fa prestiti e investimenti, la germania prende i soldi della triangolazione \t paga debiti di guerra a fr e uk, che quindi pagano i loro debiti di guerra \\
Cosi pero USA riapre una circolazione di capitali, favorevole a Stresemann \\
La germania quindi riparte \t gia dal 25:
\bi
  \item produzione supera prima guerra
  \item buona parte della pop. ha lavoro
\ei
La ripresa della germania sara ininterrotta, ma sempre dipende da stati uniti \t questo è il punto debole: con crollo di wall street, non arrivano capitali e germania in crisi \t non era autosufficiente \\
Inoltre Germania con lui normalizza i rapporti con altri paesi \t es. trattati di Locarno \t che fanno parlare dello "spirito di locarno" = clima di distensione \\
Con trattato:
\bi
  \item perdita alsazia lorena
  \item accetta la smilitarizzazione renania
  \item no accordi per confini orientali, dove infatti si espandera hitler \t la germania dovra trattare direttamente con i singoli stati
\ei
Nel 26 Germania entra anche in Societa delle Nazioni \t e firma anche trattati con Russia \\
Dal 25 quindi, la germania smette di essere considerata il paese vinto \t torna invece ad essere protagonista \\
Patto Biran-Kellog (parigi, 1928, 15 paesi) \t i paesi si impegnano a rinunciare alla guerra per via violenta, ma di risolvere conflitti per via diplomatica \t anche URSS si aggiungera (Brian ministro francese, Kellog sgretario USA) \\
Nel 25 pero ci sono elezioni del presidente \t Erbert era morto, eletto Hindenburg (che pero era monarchico) \t accade perche partiti comunisti sostengono candidati distinti, non fanno fronte unico \\
Nel 28 invece elezioni politiche \t nuovo governo di colazione (cattolici popolari liberali) con Muller \t divisioni interne, governo debole \t e rimane il problema delle riparazioni di guerra \\
I social democratici non volevano ridurre spese sociali, mentre altri nel governo si \\
Nel 28 muore Stresemann \t nel 29 crollo di Wall Street \t disastro \\
Muller si dimette nel 30, governo a Bruning (vicino a Hindenburg) \t lui è l'autore della disfatta della repubblica di Weimar \\
Tensioni da estrema sinistra e destra \t Bruning governa per decreti con appoggio di Hindenburg \t fa appello all articolo 48 (pretesto per aggirare ogni prassi democratica, che crea opposizione parlamentale) \\
Fino al 32 governa Bruning \t nel 30 inoltre nuove elezioni \t scontri tra nazisti e comunisti \t Hitler ottiene 6 milioni e mezzo di voti \\
Nel 32 elezioni presidenziali, Hitler non eletto ma ha grande successo di voti \t riletto Hindenburg, che licenzia Bruning \t chiama al governo uomini + conservatori, ma governi durano poco \\
Nel frattempo crisi economica \t si va di nuovo al governo (due volte nel 32) \t in entrambi casi il governo non esce rafforzato:
Nel luglio, nazisti diventano 1 partito germania \t rivendica gia pos. di cancelliere \t ha appoggio di industria ed esercito \t appoggio non solo politico, ma anche economico \t sono ritenuti l unica forza di poter portare all ordine \\
Nel 33 Hindenburg affida quindi a Hitler di creare nuovo governo

\ss[Nazismo]
Mein kampf:
\bi
  \item lotta a liberalismo e marxismo \t perche parlmanetarismo e democrazia sono deboli, non si puo difendere la liberta individuale ne la lotta di classe \t se no no unita dello stato
  \item lotta contro ebrei = incarnano sfruttamento economico, male della germania
  \item razza superiore, che deve espandere il suo lebensraum (verso est)
  \item combattere bolscevismo e operai
\ei
i
"Colpo alla schiena" \t Germania non ha perso la guerra, ma la condizione è tale per colpa di ebrei e marxisti \\
In Italia c'è mito vittoria mutilata, mentre in germania c'è rifiuto degli esisiti della WWI \t sconfitta non attribuita a debolezza germania, ma a cause esterne \\
Il tema del razzismo non è solo di Hitler in realta \t es. Rosenberg con "Il mito del 20esimo secolo", che è praticamente ideologo del nazismo \\
Hitler prende il potere:
\bi
\item incendia parlamento nel 33, da colpa ai comunisti
\item quindi misure eccezionali, articolo 48, sospesi diritti fond
\item Hindenburg scioglie parlamento, nuove elezioni \t nazisti vincono insieme a socialisti, che hanno maggioranza in parl
\ei
Hitler quindi va al potere in modo legale \t no colpo di stato qua: vince elezioni, va in parlamento \\
Nel marzo 33 parlamento vota legge dei pieni poteri \t vota la sua esauturazione, come è possibile? \t maggioranza è nazista \\
Chi vota contro la legge sono social democratici (comunisti erano stati cacciati) \t inizia quindi la dittatura di Hitler \\
Costruisce Gestapo, SS, magistratura sotto governo, no partiti politici, no sindacati, notte dei lunghi coltelli nel 34 elimita capi dell SA (Rommel che lo contesta) \\
Nel 34 muore Hindenburg, questo quindi segna la fine della repubblica di Weimar \t non ci sono elezioni presidenziali, lui diventa Fuhrer e inaugura il 3 Reich \\
Esodo di intellettuali \\
Hitler firma nel 33 con Pio 11 \t firma trattato con chiesa, che garantisce indipendenza della chiesa catt. in germania \t la chiesa lo fa perchè vuole mantenere un margine di controllo in germania \t con enciclica "Con cocente dolore" prende invece distanze \\
Dopo enciclica cattolici perseguitati \t chiesa protestante invece sottomessa \\
Persecuzione ebrei da subito:
\bi
  \item 33-35: discriminazione, grande propaganda per diffondere ostilita nei confronti degli ebrei, licenziamento dalle cariche pubblica etc.
  \item 35: leggi di Norimberga, che privano gli ebrei della cittadinanza = non avere diritto civile ne politico nel paese, sono dichiarati razza inferiore \t quindi esclusi da molti lavori, non si possono sposare ebrei etc.
  \item 1938: ucciso tedesco a Parigi da ebreo, quindi scusante per Notte dei cristalli
\ei
Inizio della soluzione finale inizia con invasione della URSS \t nel 1942 conferenza di Van See, in cui progettano la soluzione finale \t nel 42 inizia anche deportazione intensa

\ss[Politica economica di Hitler]
Prevede due settori: nel settore agricolo e in quello industriale \\
In quello agricolo si vuole autarchia (come in italia, ma qua in funzione della guerra) \t corporazione alimentare del Reich creata, che controlla consumo, produzione di tutti i beni (monopolio di stato) \\
Poderi fino a 125 ettari sono patrimonio dei contadini \t piccoli proprietari terrieri tutelati \t si ritorna al feudalesimo quasi, in cui il contadino è legato indissolubilmente alla sua terra \\
I latifondisti ottengono grandi sovvenzioni statali, e si arricchiscono \\
Per il settore industriale c'è massimo sforzo \t la crescita è finalizzata al riarmo \\
Nel 35 leva obbligatoria, nel 36 grandi commesse militari \t e nel 36 infatti hitler inizia a preparare la guerra \t nel 36 piano quadriennale (imita quelli di stalin) \t costruiti canali, strade etc. per riassorbire disoccupazione \\
Nel 38 disoccupazione riassorbita completamente \t grazie anche a serie di leggi tra 34 e 35, tra cui l'impedimento di scegliere il proprio lavoro (stato da lavoro in base alla necessita), e il lavoro obbligatorio per giovani tra 18 e 25 anni per avere + manodopera \\
Questo accade in un contesto senza sindacati (esiste solo "associazione dei lavoratori????" che ha carattere corporativo come in italia) \\
Queste manovre generano grande consenso \t non c'è liberta sindacale, ma al popolo interessa di + avere un lavoro \\
I salari rimangono bassi, pero c'è un controllo dei prezzi da parte dello stato \t inoltre lo stato si assume supporto per le pensioni etc. \t la percezione è che i lavoratori sono favoriti (anche con volkswagen, arma di propaganda fortissima \t la macchina del popolo)

\ss[Politica estera di Hitler pre WWII]
Nel 34 tenta di annetter austria, ma fallisce, mussolini media, "guardie del brennero" \\
Smantella il trattato di Versaille e viola i trattati internazionali:
\bi
  \item nel 35 riprende la Saar con un plebiscito
  \item nel 36 invade la Renania (in realta manda una piccola parte dell'esericito) \t francia teme vera invasione, quindi si ritira
  \item trova nuovo alleato in Italia, con guerra di spagna anche
  \item nel 37 asse roma berlino tokio ??
\ei
Europa si riunisce nella conferenza di Stresa, in cui condannato operato di Hitler ma solo a parole \t inoltre si mostra la debolezza della Societa delle Nazioni, che non interviene (esempio in Renania) \\
Cause WWII: 
\bi
  \item politica violenta di hitler
  \item politica di accondiscendenza 
\ei
Nel 38 c'è conferenza di Monaco, Hitler rivendica i Sudeti \t Cecoslovacchia non vuole cederli, ma non possono partecipare alla conferenza \t UK e FR cedono quindi \t pero hitler deve garantire che moravia e ? devono rimanere libere (cosa che non rispettera) \\
Nel 38 si prende austria, 2 tentativo di annessione funziona \t non viene presentata come un invasione ma come invito di un governo filo nazista \\
Nel 39 prendere Boemia e Moravia \t poi chiede alla polonia la città di Danzica \\
A questo punto UK e FR dicono che se hitler invade polonia, allora interverranno \\
Italia occupa albania nel ? \t e Mussolini rivencida anche Savoia, Corsica, Tunisi \t nel 39 firmato il patto d'acciaio \\
Il 23 agosto pero firma anche patto Ribbentrop-Molotov \t patto di non aggressione di 10 anni con urss \t aveva bisogno di questo perche doveva avere sicurezza del non intervento in fronte occidentale \\
Stalin accetta perche in realta è prevista la spartizione della polonia, e i baltici e scandinavia vengono indicate come zone di espansione \\
Patto inaspettato, i due si odiano \t definito il patto + cinico della storia \\
Hitler invade il 1 settembre la polonia del 1939 \\
WWII \t guerra + ideologica della WWI, perche si contrappongono due ideologie: potenze democratiche contro nazifascismo

\ss[Crisi di Wall Street]
In USA accade che nei ruggenti anni 20 \t momento di benessere e crescita in USA (mentre in europa momenti difficili) \t usa aveva capitali in eccedenza che investe in europa, con triangolazione finanziaria \\
In USA c'è quindi corsa all'arricchimento \t quindi crea una corsa agli investimenti \t c'è grande stabilita produttiva, tutti corrono a investire \\
Grande ottimismo, che convive con politica fortemente isolazionista \t dopo essere intervenuti in WWI, non entrano in Societa nazioni \\
Harding vince le presidenziali nel 20, ed è l'emblema dell'isolazionismo (ed è repubblicano) \t anche probizionismo e xenofobia \\
Programma econmico dei repubblicani è liberista \t ridotta spesa pubblica, ridotte imposte dirette, abbassato accesso al credito = si rinuncia a regolaizone economia \\
Il valore delle azioni cresce, ma poi c'è scollamento tra piano economico e finanziario \t valore in borsa continua a crescere, mentre aziende in realta si fermano \t perche c'è crisi di sovrapproduzione e saturazione del mercato = rallentamento econmico che non corrisponde a un rallentamento finanziario \\
Qualcuno ha sentore questo, e quindi vende le sue azioni \t nell ottobre del 29 crolla la bolla speculativa \t operatori di borsa liquidano i loro titoli, e tutti cominciano a fare stessa cosa \\
Il giovedi nero cedute 13 milioni di azioni \t azioni non valgono + niente, quelli che le possiedono le vogliono vendere \\
Reazione a catena: i risparmiatori non possono + ripagare debito con banca, le banche non hanno + liquidita e quindi chi aveva soldi nelle banche vanno persi (inoltre banche avevano dato prestiti ad aziende) \\
Licenziamenti di massa \t alcune chiudono, altre devono sostenere costi = grande problema sociale \t 14 milioni di disoccupati nel 1932 \\
Banche falliscono, anche perche stato non interviene praticamente \t il governo è assente perche non si prende alcuna decisione anche se si sarebbe dovuto abbassare il tasso di interesse e altri interventi \\
Hoover credeva che fosse una crisi ciclica che avviene ogni 10 anni \t e anzi approva interventi protezionistici contro parere di economisti \\
Finisce triangolazione \t quindi europa vive anche la crisi, che si ripercuote in germania \t conferenza di Losanna viene ratificata l impossibilita della germania di pagare il debito di guerra

\ss[Roosevelt]
Propone il new deal, vuole attuare politica interventista di supporto alla popolazione \t per far partire economia \\
Roosevelt passa da politica liberista a una in cui lo stato è presente \\
Bisogna rilanciare gli investimenti delle aziende e i consumi dei cittadini \t fa politica pragmatica, senza creare questoine ideologia \\
Segue il pensiero di Kanes \t stato deve essere protagonista del mercato e economia liberista non funziona

\sss[Interventi indiretti]
Provvedimento che migliorano le condizioni delle attivita produttive e popolazione \\
Dollaro viene sganciato dall'oro, la moneta viene svalutata \\
Si apre mercato estero per risolvere sovrapproduzione e incentivate esportazioni \\
Legge agricola: da premi in denaro a agricoltori che limitano la produzione \t si vuole limitare la sovrapproduzione nel settore agricolo \\
Introdotto un codice di disciplina \t aziende fanno accordi vincolanti per limitare produzione \\
Inoltre si impone la rinuncia al lavoro nero, al lavoro minorile, orario di lavoro stabilito \\
Riforma fsicale, in cui si ha tassazione progressiva \t aliquote + alte per i redditi + alti \\
Wagner act \t sancisce diritto di sciopero e contrattazione collettiva (contratto dei metalmecc, degli insegnanti etc.)

\sss[Interventi diretti]
Lo stato interviene pero anche direttamente, per riassorbire sovrapproduzione \t crea quindi cantieri pubblici \\
Costruite dighe nel bacino del Tennessee \t si vuole sfruttare energia idroelettrica \\
Si costruiscono strade, canali, ponti, etc. \t iniziative danno lavoro, e mettono a disposizione dell industria e dell agricoltura prezzi bassi \\
INoltre nuovo sistema assistenziale, che prevede sussidi di disoccupazione e assistenza \t finanziato dalle tasse e dalla casse dello stato

\v

USA esce dalla crisi lentamente, ma new deal ha successo \t Roosevelt ha carisma, trasmissione radiofonica in cui aggiorna popolazione sui progressi \\
Non piace a tutti \t va a stravolgere il sistema liberista \t trova opposizione nella corte suprema, che rappresenta anche grandi imprenditori \\
Lui viene accusato di creare politica incentrata sulla sua persona \t condannata la troppa ingerenza dello stato, e del presidente \\
Roosevelt vince questo scontro, perche nel 37 riesce a sostituire alcuni giudici con elementi + favorevoli \\
Il new deal modifica il rapporto tra politica e societa \t getta le basi del "wealthfare state" (assistenza ai popoli, pensioni, etc. assunte dallo stato) \\
New deal espande pero anche burocrazia, perche crea organismi di stato \\
Inoltre nuovo rapporto tra politica ed economia \t stato deve regolare il sistema economico, ma non intervenire pesantemente \\
Inoltre cambia percezione del ruolo dei sindacati \t che non sono degli elementi da combattere, ma degli interlocutori importanti per dare via istituzionale ai problemi sociali \\
Sindacati possono anche portare alle istituzioni i problemi sociali

\ss[Spagna]
Spagna è un paese arretrato, che vive in una condizione di autoritarismo monarchico \t ha economia basata sul latifondo \\
Presenza di sindacati anarchici e socialisti \t vogliono riforme agrarie, e spinete autonistiche in paesi baschi e catalogna \\
Alfonso 13 (borbone) favorisce colpo di stato di Primo De Rivera \t si ispira al fascisto, rimane in carica fino al 1930 e alfonso rimane il sovrano \\
Nel 31 alle amministrative vincono paritti repubblicani, re va in esilio procalamata repubblica \t ma si alternano forze di sinistra e di destra (rappresentate dall falange, il partito fascista spagnolo) \\
Falange era stato fondato da Antonio, figlio di Primo \t partito si ispira al fascismo italiano direttamente, ma non e un movimento di massa \\
Nel 1936 il fronte popolare (coalizione di tutte le forze di sinistra, con anche liberali e democratici) vincono elezioni \t si uniscono forze antifasciste \\
Il fronte popolare vede anche la presenza dei comunisti: stalin da direzione per impedire nuovo governo fascista \\
Governo è sostenuto esternamente dai socialisti \t ma govenro non fa riforme che aveva promesse (la coalzione è troppo eterogenea) \t quindi rivolte violentissime contro il clero e contro i proprietari \t tensione sociali \\
Quindi destra fa il colpo di stato \t in realta la destra non aveva riconosciuto il risultato delle elzioni e non riconosceva il governo di sinistra ????? \\
Francisco franco con l aviazione tedesca e truppe italiane, sbarca dal marocco nel sud della spagna \t guerra civile \\
Vince nel 39 quando cadono madrid e barcellona \\
Nessuno aiuta il fronte popolare (europa non si muove) \t infatti francia e ingh. e italia e germania firmano patto per non intervento \\
Francia infatti non voleva tensioni con ingh., mentre ingh. non voleva deriva sinistra \t mentre urss attraverso il commintern interviene, che manda aiuti e crea le brigate (associazioni di volontari) \t ma sono inferiori \\
Esercito repubblicano è anche frammentato dalle divisioni interne, che si scontrano \\
Francisco instaura dittatura, che durerara fino al 75 \\
Guerra spagna è importante perche avvicina italia e germania \\
Hitler prova qua mezzi che non erano mai stati provati \t testa gli armamenti per la WWII \\
Opus Dei \t associazione il cui fondatore ha simpatie franchiste prob.
\end{document}
